\begin{abstract}
\textbf{\textit{Contribution:} This research identifies the critical, often-overlooked role of active verbal narration in hands-on instructional videos, demonstrating that audio transcript quality directly impacts semantic learning alignment and viewer engagement metrics. \textit{Background:} While robotics education traditionally relies on high-cognitive-load, lecture-style videos, a growing volume of online instructional content utilizes silent, demonstration-only videos. However, the exact empirical impact of missing instructor audio narration on learner engagement remains poorly understood. \textit{Research Questions:} 1) Is audio a significant component of a video designed to teach robotics? 2) To what extent can embodied learning be leveraged through only using audio in robotics education? \textit{Methodology:} A structured pipeline utilizing automated speech recognition (ASR) and Large Language Models (LLMs) was developed to analyze 130 robotics instructional videos categorized into conventional and embodied modalities. Because embodied instructional videos on public platforms are produced in two fundamentally distinct narration styles, the study design stratified the embodied category a priori by speech rate into verbally narrated tutorials ($n=48$) and silent demonstrations ($n=26$) prior to the comparative analysis. Quantifiable metrics ($S_{corr}$, speech rates, phrase counts) were scored against reflection datasets from a dual-cohort user study consisting of 65 undergraduate engineering technology students at Middle Tennessee State University and broader, spontaneous YouTube audience comments, complemented by a post-viewing video preference survey of the same 65 participants. The study is limited by its primary reliance on transcript-based text analysis rather than direct multi-modal visual frame evaluations. \textit{Findings:} The speech rate analysis confirmed the designed stratification, exhibiting a clearly bimodal distribution in which silent demonstration videos generated hallucinated or negligible transcripts; a sensitivity analysis showed that aggregating the two embodied subgroups would systematically understate embodied performance. Within the planned three-group comparison, verbally narrated embodied tutorials performed on par with or exceeded conventional lectures across nearly all semantic metrics. Verbally embodied instructors achieved higher keypoint reinforcement scores (3.67 vs. 3.58) and a 20\% higher density of physical action language, driving substantially higher public engagement rates on YouTube. The preference survey corroborated these findings: hands-on demonstration (88\%) and clear verbal explanation (85\%) were the two most frequently selected instructional features, whereas body language alone was selected by only 28\% of participants, underscoring that narrated action rather than visible gesture carries the instructional value of embodied videos.}
\end{abstract}